\documentclass{article}
\usepackage{mathpartir}

\begin{document}

% infix:53 " ∈ " => "\in"
% infix:53 " ∉ " => "\notin"
% infix:53 " ≠ " => "\ne"
% infix:53 " = " => "="
% infixr:35 " ∧ " => "\land"
% infix:50 " → " => "\mathrel{\to_{\beta}}"
% infix:50 " ↠ " => "\mathrel{\to_{\beta}^{*}}"
% notation:56 "λ" x:atom "." e:56 => "\lambda " x ".\," e
% notation:75 "FV" e:76 => "\mathrm{FV}\," e
% notation:53 z:atom "fresh" => z "\;\mathsf{fresh}"
% notation:85 e:85 "[" s:0 "/" x:atom "]" => e "[" s "/" x "]"

Variables range over names, and terms are identified up to alpha-equivalence.

\begin{figure}
\[
\begin{array}{l}
\mbox{Terms} \\
\begin{array}{rcl}
e,f,s & ::= & x \mid e_{1}\,e_{2} \\
 & & \lambda x.\,e
\end{array}
\end{array}
\]
\caption{Untyped lambda syntax}
\end{figure}

Substitution replaces free occurrences and avoids capture. In the guards,
\(\mathrm{FV}\,s\) denotes the free variables of \(s\). In the final equation,
\(z\;\mathsf{fresh}\) means that \(z\) occurs nowhere in \(e\) or \(s\) and
is distinct from \(x\) and \(y\).

\begin{figure}
\[
\begin{array}{rcll}
x[s/y] & = & s & \mbox{if}\;x=y \\
x[s/y] & = & x & \mbox{if}\;x\ne y \\
{} & {} & {} & {} \\
\left(e_{1}\,e_{2}\right)[s/x] & = & \left(e_{1}[s/x]\right)\,\left(e_{2}[s/x]\right) &  \\
{} & {} & {} & {} \\
\left(\lambda x.\,e\right)[s/x] & = & \lambda x.\,e &  \\
\left(\lambda y.\,e\right)[s/x] & = & \lambda y.\,e[s/x] & \mbox{if}\;y\ne x\land y\notin\mathrm{FV}\,s \\
\left(\lambda y.\,e\right)[s/x] & = & \lambda z.\,e[z/y][s/x] & \mbox{if}\;y\ne x\land y\in\mathrm{FV}\,s\land z\;\mathsf{fresh}
\end{array}
\]
\caption{Capture-avoiding substitution}
\end{figure}

\begin{figure}
\begin{mathpar}
\inferrule*[right={R-Beta}]
  { }
  {\left(\lambda x.\,e\right)\,s\mathrel{\to_{\beta}}e[s/x]}

\inferrule*[right={R-AppLeft}]
  {e\mathrel{\to_{\beta}}f}
  {e\,s\mathrel{\to_{\beta}}f\,s}

\inferrule*[right={R-AppRight}]
  {e\mathrel{\to_{\beta}}f}
  {s\,e\mathrel{\to_{\beta}}s\,f}

\inferrule*[right={R-UnderBinder}]
  {e\mathrel{\to_{\beta}}f}
  {\lambda x.\,e\mathrel{\to_{\beta}}\lambda x.\,f}
\end{mathpar}
\caption{Beta reduction}
\end{figure}

\begin{figure}
\begin{mathpar}
\inferrule*[right={R-ManyRefl}]
  { }
  {e\mathrel{\to_{\beta}^{*}}e}

\inferrule*[right={R-ManyStep}]
  {e\mathrel{\to_{\beta}}f \\
   f\mathrel{\to_{\beta}^{*}}g}
  {e\mathrel{\to_{\beta}^{*}}g}
\end{mathpar}
\caption{Multi-step reduction}
\end{figure}

\end{document}
